\documentclass[10pt,a4paper]{extarticle}
\usepackage{parskip}

\usepackage[utf8]{inputenc}
\usepackage[T1]{fontenc}
\usepackage[brazil]{babel}
\usepackage{amsmath, amssymb, amsthm}
\usepackage{geometry}
\usepackage{xcolor}
\usepackage{hyperref}
\usepackage{booktabs}
\usepackage{array}
\usepackage{tabularx}
\usepackage{enumitem}
\usepackage[expansion=false,protrusion=false]{microtype}
\usepackage{csquotes}

\definecolor{important}{RGB}{255, 100, 100}
\definecolor{notice}{RGB}{0, 102, 204}

\newcommand{\impt}[1]{\textcolor{important}{#1}}
\newcommand{\ntc}[1]{\textcolor{notice}{#1}}

\setlist{nosep}
\geometry{margin=2.2cm}

\hypersetup{
    colorlinks=true,
    linkcolor=notice,
    urlcolor=notice,
    citecolor=notice
}

\title{Estrutura Acadêmica da UFOP: Departamento e Colegiado de Curso}
\author{Introdução ao Curso -- DECOM/UFOP}
\date{\today}

\begin{document}
\maketitle

\begin{abstract}
\noindent Estas notas apresentam dois órgãos centrais da vida acadêmica do estudante na UFOP: o \impt{departamento} (ou organização de nível hierárquico equivalente) e o \impt{colegiado de curso}. Entender quem decide o quê é fundamental para o estudante saber a qual instância recorrer em cada situação --- seja para pedir o aproveitamento de uma disciplina cursada em outra instituição, seja para sugerir alteração em uma ementa, seja para participar da representação discente. O texto se apoia diretamente no \emph{Estatuto da UFOP} (Resolução CUNI nº 1.868/2017) e no \emph{Regimento Geral} (Resolução CUNI nº 1.959/2017).
\end{abstract}

\tableofcontents
\vspace{0.5cm}

\section{Onde estamos na estrutura da Universidade}

Antes de entrar nos dois órgãos, vale uma localização rápida. A UFOP se organiza, no nível mais alto, em \impt{Administração Central} (Reitoria, conselhos superiores, Conselho Universitário, Conselho Curador) e \impt{Unidades Acadêmicas}. Uma unidade acadêmica é, por exemplo, o ICEB (Instituto de Ciências Exatas e Biológicas), ao qual o DECOM está vinculado.

Cada unidade acadêmica é, por sua vez, composta por:

\begin{itemize}
    \item conselho de unidade acadêmica;
    \item diretoria e vice-diretoria;
    \item cursos de graduação e pós-graduação;
    \item \impt{colegiados de cursos} de graduação e de pós-graduação;
    \item \impt{departamentos} (ou organizações de nível hierárquico equivalente);
    \item núcleos e órgãos complementares.
\end{itemize}

Os dois últimos itens em destaque são o assunto desta aula. A distinção essencial é a seguinte:

\begin{itemize}
    \item o \impt{departamento} é a unidade de \emph{lotação} dos docentes --- é \emph{onde o professor trabalha};
    \item o \impt{colegiado de curso} é a instância de coordenação \emph{didática} de um curso específico --- é \emph{quem cuida da formação do estudante daquele curso}.
\end{itemize}

\ntc{Um docente do DECOM, por exemplo, pode ministrar disciplinas para alunos do BCC, da Engenharia de Computação, da Estatística e da Matemática. Ele é lotado no DECOM (departamento), mas eventualmente é representante de departamento em mais de um colegiado de curso.}

\section{Departamento}

\paragraph{Definição.} O departamento (ou organização de nível hierárquico equivalente) é, conforme o Estatuto da UFOP, o \impt{órgão de lotação de professores} e \emph{um} dos órgãos de lotação dos técnicos-administrativos em educação, voltado para \emph{objetivos comuns de ensino, pesquisa e extensão}, em uma determinada área do conhecimento. Toda Unidade Acadêmica deve ser composta obrigatoriamente por departamentos (ou estruturas equivalentes), e todo departamento existente está vinculado a alguma unidade acadêmica.

\paragraph{Por que ele existe.} A existência de um departamento se justifica pelas \emph{áreas do conhecimento que ele abrange}, pelas \emph{linhas de pesquisa e projetos pedagógicos} que sustenta, e pelos \emph{recursos materiais e humanos} necessários ao seu funcionamento. O DECOM existe porque há um conjunto coerente de áreas (algoritmos, IA, sistemas, teoria da computação, etc.), pesquisa associada, e pessoal alocado para sustentar tudo isso.

\paragraph{Quem o compõe.} A \impt{Assembleia do Departamento} é o órgão deliberativo do departamento. Ela reúne:

\begin{itemize}
    \item \impt{todos os docentes} lotados no departamento;
    \item representante(s) dos \impt{técnicos-administrativos em educação} lotados no departamento;
    \item representante(s) do \impt{corpo discente} dos cursos atendidos pelo departamento.
\end{itemize}

O número de representantes discentes e técnico-administrativos é paritário entre si. As deliberações são tomadas por maioria dos presentes.

\paragraph{Chefia e mandato.} A assembleia escolhe o \impt{chefe} e o \impt{vice-chefe} do departamento, com mandato de \impt{2 anos}, permitida uma recondução. O chefe preside a assembleia.

\paragraph{Câmara departamental.} Departamentos com \emph{mais de 20 docentes} podem criar uma \impt{Câmara} para deliberar sobre os assuntos do departamento de forma mais ágil. A câmara, quando existe, é composta pelo chefe, vice-chefe, cinco representantes docentes, um técnico-administrativo e um discente. Nesse caso, a assembleia plena continua existindo, mas é convocada com periodicidade reduzida (ao menos uma vez por semestre).

\paragraph{Reuniões.} Os departamentos devem se reunir em assembleia \impt{pelo menos duas vezes por semestre}. Quando há câmara, é exigida \emph{pelo menos uma} assembleia plena semestral.

\subsection{O que o departamento resolve}

A pergunta operacional para o estudante é: \emph{quando preciso interagir com o departamento?} O departamento decide sobre tudo aquilo que envolve o \emph{trabalho dos docentes} naquela área de conhecimento. Concretamente, compete à Assembleia do Departamento:

\begin{enumerate}[label=(\arabic*)]
    \item \impt{Planejar as atividades de ensino, pesquisa e extensão} do departamento --- definir o que será ofertado, como, e por quem.
    \item \impt{Elaborar os planos de trabalho e de capacitação} dos docentes e técnicos lotados nele. Isto inclui afastamentos para doutorado, pós-doutorado, capacitação, etc.
    \item \impt{Atribuir encargos de ensino, pesquisa e extensão} aos docentes. É o departamento que decide \emph{quem dá qual disciplina} em cada semestre, respeitando o planejamento e as linhas de pesquisa.
    \item \impt{Propor aos colegiados de curso} os programas, ementas, cargas horárias, pré-requisitos e bibliografias das \emph{disciplinas que o departamento oferece}.\\
    \ntc{Atenção à direção: o departamento \emph{propõe} a ementa de PCC104 (Projeto e Análise de Algoritmos), por exemplo; cabe ao colegiado do BCC e demais colegiados que utilizam a disciplina aprovar essa ementa no contexto do projeto pedagógico de cada curso.}
    \item \impt{Propor ao Conselho da Unidade} a contratação, substituição, afastamento e dispensa de docentes.
    \item \impt{Eleger os representantes do departamento nos colegiados de curso}. Como o DECOM atende a múltiplos cursos, ele indica representantes em diversos colegiados.
    \item \impt{Elaborar o relatório anual de atividades} do departamento, com base nos relatórios individuais dos docentes. Esse relatório é homologado pelo Conselho da Unidade.
    \item Propor, com 2/3 dos membros, o \impt{afastamento ou destituição do chefe}.
\end{enumerate}

\paragraph{Resumo operacional.} O departamento responde pela \impt{oferta} (quem ensina, o que ensina, com qual conteúdo de disciplina) e pelo \impt{corpo docente} (contratação, encargos, capacitação). Pesquisa e extensão são planejadas e organizadas no departamento. Se você quer entender por que uma disciplina tem determinada ementa, ou por que um certo professor está ministrando uma determinada matéria, a resposta começa no departamento.

\subsection{Chefia do departamento: atribuições}

O chefe do departamento, escolhido pela própria assembleia, é o representante executivo do órgão. Compete a ele:

\begin{itemize}
    \item representar o departamento;
    \item coordenar, no plano executivo, as atividades de ensino, pesquisa e extensão;
    \item apresentar o relatório anual de atividades ao Diretor da Unidade, após apreciação pela Assembleia;
    \item cumprir e fazer cumprir as deliberações da Assembleia e dos órgãos superiores;
    \item controlar a frequência dos docentes e técnicos, bem como a \emph{execução dos planos de ensino};
    \item exercer o poder disciplinar em sua esfera de jurisdição.
\end{itemize}

\section{Colegiado de Curso}

\paragraph{Definição.} O colegiado de curso é o órgão responsável pela \impt{coordenação didática} dos componentes curriculares do projeto pedagógico de um curso específico. Cada curso de graduação e cada programa de pós-graduação tem o seu próprio colegiado.

\ntc{Note a mudança de foco em relação ao departamento. O departamento se organiza por área de conhecimento e abriga docentes; o colegiado de curso se organiza por curso e cuida da trajetória formativa do aluno daquele curso. Um aluno do BCC tem um colegiado --- o do BCC --- mesmo que cursando disciplinas de cinco departamentos diferentes.}

\paragraph{Quem o compõe.} O colegiado é constituído por:

\begin{itemize}
    \item \impt{representantes docentes} (e/ou técnico-administrativos diretamente envolvidos em atividades didáticas) dos departamentos que oferecem disciplinas do curso;
    \item \impt{representação estudantil}, eleita pelos pares, com mandato de 1 ano (permitida uma recondução).
\end{itemize}

A proporção de representação docente segue a regra do Regimento Geral: \impt{1 representante a cada 180 horas} de carga horária que o departamento oferece para o curso. Departamentos com carga inferior a 180h podem participar (faculdade); a representação é limitada a, no máximo, \impt{5 membros por departamento}.

\paragraph{Coordenação e mandato.} O colegiado elege, entre seus membros, um \impt{coordenador de curso} e um \impt{vice-coordenador}, ambos docentes, para mandato de \impt{2 anos} (permitida uma recondução). O coordenador preside o colegiado; na ausência dele, o vice-coordenador o substitui.

\paragraph{Reuniões.} Os colegiados devem se reunir \impt{ao menos duas vezes por semestre}. Devem também elaborar periodicamente relatório de acompanhamento dos processos de ensino-aprendizagem e da \impt{integralização curricular}.

\paragraph{Núcleo Docente Estruturante (NDE).} No caso de cursos de graduação, o colegiado indica os membros do \impt{NDE} ou órgão similar --- um grupo menor, focado em pensar e atualizar o projeto pedagógico do curso. Os indicados podem ser ou não membros do colegiado.

\subsection{O que o colegiado resolve}

O colegiado é o órgão que, do ponto de vista do estudante, decide \emph{o que conta no seu histórico}. Compete a ele:

\begin{enumerate}[label=(\arabic*)]
    \item \impt{Compatibilizar as diretrizes gerais dos componentes curriculares} do curso e propor modificações.
    \item \impt{Regulamentar os componentes curriculares} para execução do projeto pedagógico.
    \item \impt{Deliberar sobre ementas e programas} elaborados pelos departamentos, no contexto do projeto pedagógico do curso. \ntc{Ou seja: o departamento propõe; o colegiado aprova (ou pede ajuste) à luz do projeto pedagógico do curso.}
    \item \impt{Propor ao Conselho Superior de Graduação (ou Pós-Graduação)} o projeto pedagógico do curso e suas alterações --- incluindo pré-requisitos, carga horária, ementas, programas, regulamentos e componentes curriculares.
    \item Decidir sobre questões individuais do percurso acadêmico do aluno:
    \begin{itemize}
        \item \impt{reopção} de cursos (mudança interna de curso);
        \item \impt{equivalência de disciplinas} (uma disciplina cursada substitui outra);
        \item \impt{desligamento e jubilamento};
        \item \impt{aproveitamento de estudos};
        \item ingresso de \impt{portador de diploma} de graduação em novo curso;
        \item \impt{transferência} de outras instituições;
        \item \impt{reingresso};
        \item \impt{mobilidade acadêmica} nacional e internacional.
    \end{itemize}
    \ntc{Esse conjunto inclui basicamente toda decisão que envolve ``minha matéria $X$ feita na situação $Y$ vale para o meu curso?''. A resposta é dada pelo colegiado, não pelo departamento que oferece a disciplina.}
    \item Apreciar recomendações da unidade acadêmica e requerimentos de docentes sobre assuntos do curso.
    \item \impt{Coordenar a orientação acadêmica} dos estudantes, visando à integralização curricular e à colação de grau.
    \item \impt{Indicar à Pró-Reitoria competente os candidatos à colação de grau} ou à diplomação. É o colegiado que diz oficialmente: ``este aluno cumpriu tudo o que era exigido pelo projeto pedagógico''.
    \item Indicar membros do NDE (no caso de graduação).
    \item Recomendar ao departamento responsável providências relativas à \impt{abertura de vagas e de turmas} e à melhor utilização de instalações, materiais e pessoal.
\end{enumerate}

\subsection{Resumo operacional}

O colegiado responde por \impt{trajetória}, \impt{currículo} e \impt{validação de percurso}: o que o curso é, como ele se estrutura, e como o histórico individual do aluno se encaixa nessa estrutura.

\section{Departamento \emph{vs.} Colegiado: tabela comparativa}

A diferença operacional pode ser resumida da seguinte forma:

\begin{center}
\renewcommand{\arraystretch}{1.35}
\begin{tabularx}{\textwidth}{@{} l X X @{}}
\toprule
\textbf{Dimensão} & \textbf{Departamento} & \textbf{Colegiado de Curso} \\
\midrule
Organiza-se por & área de conhecimento & curso específico \\
Lotação de & docentes e técnicos & ningu\'em é \emph{lotado} no colegiado \\
Atribui & encargos didáticos a docentes & componentes curriculares ao projeto pedagógico \\
Propõe & ementas e programas das disciplinas que oferece & projeto pedagógico do curso \\
Aprova & planos de trabalho de docentes & equivalências, aproveitamentos, transferências do aluno \\
Decide sobre & contratação, substituição, encargo & desligamento e jubilamento do aluno \\
Presidência & chefe de departamento (mandato 2 anos) & coordenador de curso (mandato 2 anos) \\
Discentes & representação na assembleia & representação no colegiado (1 ano de mandato) \\
Reuniões mínimas & 2 por semestre & 2 por semestre \\
\bottomrule
\end{tabularx}
\end{center}

\section{Requerimentos do aluno: a quem se dirigem?}

A PROGRAD organiza, em sua página de orientações gerais, uma lista dos principais requerimentos que o aluno faz ao longo da graduação. Esta seção sintetiza essa lista, explicitando \emph{a qual instância cada requerimento se dirige}. Quase todos são protocolizados pelo portal \texttt{minhaUFOP} (\impt{PROTOCOLIZAÇÃO DE REQUERIMENTO}), dentro dos prazos do calendário acadêmico.

\subsection{Requerimentos ao departamento}

A PROGRAD lista apenas \impt{dois tipos} de requerimento dirigidos diretamente ao departamento, ambos relacionados à oferta de disciplinas:

\begin{enumerate}[label=(\arabic*)]
    \item \impt{Abertura de turma} --- solicitação para que uma nova turma de uma disciplina seja ofertada no semestre.
    \item \impt{Abertura de vaga em disciplina} --- solicitação de vaga adicional em turma já existente, quando esta está lotada.
\end{enumerate}

\ntc{A lógica é direta: o departamento é quem oferece a disciplina --- decide encargo docente, capacidade e número de turmas. Qualquer pedido que envolva \emph{disponibilizar a disciplina} (mais turmas, mais vagas) é dele.}

\subsection{Requerimentos ao colegiado de curso}

A maior parte dos requerimentos acadêmicos do aluno vai ao colegiado. Podemos agrupá-los em cinco categorias.

\paragraph{(A) Requerimentos de matrícula em situação excepcional.} São os explicitamente listados pela PROGRAD como ``requerimentos de matrícula aos colegiados de curso'':

\begin{enumerate}[label=(\arabic*)]
    \item \impt{Quebra de pré-requisito} --- matrícula em disciplina sem ter cursado seu pré-requisito formal.
    \item \impt{Quebra de período consecutivo} --- matrícula sem respeitar o pré-requisito por bloco/período.
    \item \impt{Matrícula com excesso de horas} --- ultrapassar o limite de carga horária semestral.
    \item \impt{Matrícula com excesso de horas em disciplina facultativa} --- ultrapassar o limite, especificamente por incluir disciplina facultativa.
\end{enumerate}

\ntc{Atenção: nesses quatro casos, \impt{não é permitido o trancamento parcial} da disciplina cuja matrícula foi feita por requerimento.}

\paragraph{(B) Validação de estudos já realizados.}

\begin{enumerate}[label=(\arabic*), resume]
    \item \impt{Aproveitamento de estudos} --- aproveitamento de disciplina cursada com aprovação em outro curso ou em outra IES, mediante compatibilidade de carga horária e conteúdo programático. Exige histórico escolar e programa analítico da disciplina (Resolução CEPE 7.325).
    \item \impt{Extraordinário aproveitamento} --- dispensa de disciplina por avaliação especial, via banca examinadora. Requisitos: ter cursado pelo menos 2 semestres na UFOP, CRE $\geq 7{,}0$ e estar entre os 12\% melhores coeficientes do curso (Resolução CEPE 1.986).
    \item \impt{Exame de nivelamento em língua estrangeira} --- dispensa de disciplina de língua estrangeira por avaliação que abrange o conteúdo programático (Resolução CEPE 4.672).
\end{enumerate}

\paragraph{(C) Mudanças de vínculo e trajetória.}

\begin{enumerate}[label=(\arabic*), resume]
    \item \impt{Reopção de curso} --- transferência interna para outro curso de graduação da UFOP no mesmo agrupamento de áreas.
    \item \impt{Reopção de habilitação} --- mudança para outra habilitação do mesmo curso.
    \item \impt{Reingresso} --- retorno de aluno previamente desligado ao mesmo curso.
    \item \impt{Transferência} --- entrada de aluno vindo de outra IES brasileira em curso correspondente ou afim.
    \item \impt{Portador de Diploma de Graduação (PDG)} --- início de novo curso de graduação por quem já é graduado.
    \item \impt{Mobilidade acadêmica} (ANDIFES, nacional; convênios internacionais) --- gerida pela PROGRAD, mas requer parecer/aprovação do colegiado.
\end{enumerate}

\ntc{Os procedimentos para reopção, reingresso, transferência e PDG ocorrem por meio de \impt{processo seletivo específico} organizado pela PROGRAD (vagas residuais), mas a decisão acadêmica final --- inclusive equivalência das disciplinas já cursadas --- compete ao colegiado do curso de destino.}

\paragraph{(D) Suspensão temporária do curso.}

\begin{enumerate}[label=(\arabic*), resume]
    \item \impt{Afastamento especial} --- suspensão da matrícula e das obrigações acadêmicas, em situações especiais comprovadas, por até 4 anos. \emph{Pode ser requerido uma única vez} (Resolução CEPE 1.744, art. 18).
    \item \impt{Trancamento de matrícula} (parcial ou total) --- dentro do prazo do calendário acadêmico é feito direto pelo \texttt{minhaUFOP}; \emph{fora do prazo}, depende de motivo de força maior e aprovação do colegiado (Resolução CEPE 1.744, art. 19).
    \item \impt{Licença maternidade/paternidade} (natural ou por adoção) --- trancamento total ou parcial \emph{sem contar} para o prazo máximo de integralização nem para o limite de trancamentos (Resolução CONGRAD 250/2025, art. 9).
    \item \impt{Regime Especial de Trabalho Escolar e Frequência (RETEF)} --- regime especial regulado pela Resolução CONGRAD 250.
\end{enumerate}

\paragraph{(E) Atividades curriculares e conclusão.}

\begin{enumerate}[label=(\arabic*), resume]
    \item \impt{Atividades Acadêmico-Científico-Culturais (AACC)} --- registro e validação seguem regulamento de cada colegiado, no contexto do projeto pedagógico.
    \item \impt{Segunda chamada de prova} --- aplicável apenas em caso de falecimento na família, com protocolo em até 7 dias (Resolução CEPE 2.871).
    \item \impt{Colação de grau} --- requerida pelo aluno via \texttt{minhaUFOP}, mas a \impt{indicação oficial} à PROGRAD dos concluintes é feita pelo colegiado (Art. 49, VIII do Estatuto; Resolução CEPE 5.709).
\end{enumerate}

\subsection{Requerimentos que \emph{não} vão a departamento nem a colegiado}

Importante: nem toda solicitação do aluno passa por uma dessas duas instâncias. Vão a outros órgãos:

\begin{itemize}
    \item \impt{Renovação e ajuste de matrícula}, \impt{atualização de cadastro}, \impt{emissão de documentos} (atestado de matrícula, histórico): \texttt{minhaUFOP} ou Seção de Ensino.
    \item \impt{Cancelamento de matrícula institucional} (encerramento de vínculo): PROGRAD / Seção de Ensino.
    \item \impt{Carteira estudantil}: SISBIN (Sistema de Bibliotecas).
    \item \impt{Bolsas} (alimentação, moradia, monitoria, iniciação científica, extensão, pró-ativa): pró-reitorias correspondentes (PRACE, PROGRAD, PROPP, PROEx).
    \item \impt{Auxílio à participação em eventos}: PROGRAD (Comitê de Atividades Acadêmicas).
    \item \impt{Recurso contra decisão} de departamento, colegiado ou outro órgão: instância imediatamente superior, conforme o Título VI do Regimento Geral. Tipicamente, recursos contra colegiado ou departamento vão ao \emph{Conselho da Unidade Acadêmica}.
\end{itemize}

\subsection{Tabela-resumo dos requerimentos}

\begin{center}
\renewcommand{\arraystretch}{1.25}
\begin{tabularx}{\textwidth}{@{} l X @{}}
\toprule
\textbf{Vai para...} & \textbf{Requerimentos típicos} \\
\midrule
\textbf{Departamento} & Abertura de turma; abertura de vaga em disciplina. \\
\midrule
\textbf{Colegiado de curso} & Quebra de pré-requisito; quebra de período consecutivo; matrícula com excesso de horas; aproveitamento de estudos; extraordinário aproveitamento; nivelamento em língua estrangeira; reopção (curso e habilitação); reingresso; transferência; PDG; mobilidade; afastamento especial; trancamento fora de prazo; licença maternidade/paternidade; RETEF; AACC; segunda chamada de prova; colação de grau. \\
\midrule
\textbf{Seção de Ensino / \texttt{minhaUFOP}} & Renovação e ajuste de matrícula; atualização de cadastro; emissão de documentos; cancelamento institucional. \\
\midrule
\textbf{Pró-reitorias e outros} & Bolsas (PRACE, PROGRAD, PROPP, PROEx); auxílio a eventos; carteira estudantil (SISBIN). \\
\midrule
\textbf{Conselho da Unidade} & Recursos contra decisões de departamento ou colegiado. \\
\bottomrule
\end{tabularx}
\end{center}

\paragraph{Heurística prática para o aluno.} Se o requerimento envolve \emph{disponibilizar uma disciplina} (oferta), vai ao \impt{departamento}. Se envolve \emph{qualquer outra coisa acadêmica} --- regra de matrícula, validação, mudança de vínculo, suspensão, conclusão --- vai ao \impt{colegiado}. Se é \emph{administrativo} (documento, cadastro, vínculo institucional), vai à Seção de Ensino. Se é \emph{financeiro/assistencial} (bolsa, auxílio), vai à pró-reitoria correspondente.

\section{A quem recorrer? Casos concretos}

Para fixar a distinção, considere os exemplos a seguir. A pergunta, em cada caso, é: \emph{a quem o aluno (ou o docente) deve se dirigir?}

\paragraph{Caso 1.} O aluno cursou Cálculo I em outra universidade e quer aproveitá-la no BCC da UFOP.\\
\impt{Colegiado do BCC.} É decisão sobre equivalência/aproveitamento de estudos no percurso do aluno.

\paragraph{Caso 2.} Um grupo de alunos acha que a ementa de PCC104 está desatualizada e quer sugerir uma reforma.\\
\impt{Caminho: colegiado primeiro, departamento depois.} A iniciativa de revisão pode ser levada pela representação discente ao colegiado, que dialoga com o DECOM (departamento que oferece a disciplina). A ementa, formalmente, é \emph{proposta} pelo departamento e \emph{aprovada} pelo colegiado.

\paragraph{Caso 3.} Um aluno do BCC quer mudar para Engenharia de Computação (reopção).\\
\impt{Colegiado do BCC e Colegiado da Engenharia de Computação.} Decisões sobre reopção competem aos colegiados envolvidos.

\paragraph{Caso 4.} Um docente do DECOM quer afastar-se por um ano para pós-doutorado.\\
\impt{Departamento (Assembleia do DECOM)}, que elabora os planos de capacitação e propõe afastamentos ao Conselho da Unidade.

\paragraph{Caso 5.} O DECOM precisa decidir quem vai ministrar PCC104 no próximo semestre.\\
\impt{Departamento.} A atribuição de encargos didáticos aos docentes lotados é competência da assembleia departamental.

\paragraph{Caso 6.} Um aluno concluiu todas as disciplinas e quer colar grau.\\
\impt{Colegiado do curso}, que indica oficialmente o candidato à colação de grau à Pró-Reitoria de Graduação.

\paragraph{Caso 7.} Um aluno quer questionar uma decisão tomada pelo coordenador do curso.\\
\impt{Conselho da Unidade Acadêmica}, que aprecia recursos contra atos dos coordenadores de curso (instância imediatamente superior).

\section{Por que isso importa para você}

A separação entre \emph{quem oferece a disciplina} (departamento) e \emph{quem coordena o curso} (colegiado) é uma característica estrutural das universidades brasileiras e tem uma razão clara: garante que um docente especializado em uma área do conhecimento (digamos, otimização combinatória) seja gerido em conjunto com seus pares dessa área --- mas que o projeto formativo de cada curso (BCC, Engenharia de Computação, Estatística, Matemática) seja pensado a partir da perspectiva específica do curso, integrando contribuições de vários departamentos.

Para o estudante, isso significa que há \impt{duas portas de entrada} institucionais distintas, dependendo da natureza da questão. Saber distinguir as duas evita perder tempo abrindo processos na instância errada --- e, principalmente, permite que a representação discente atue de forma estratégica em cada uma delas.

\paragraph{Representação discente.} Vale destacar que o corpo discente tem representação \emph{nas duas instâncias}: na assembleia do departamento e no colegiado de curso. São canais complementares: o representante na assembleia do departamento participa das decisões sobre oferta e corpo docente; o representante no colegiado participa das decisões sobre currículo e trajetória dos alunos.

\section{Como resolver problemas na UFOP?}

\begin{enumerate}
    \item Consultar um colega de turma - É quem está na situação mais parecida com a sua. 
    \item Consultar um veterano - Poder ter passado por situação semelhante a sua recentemente.
    \item Consultar a página de orientações \url{https://www.prograd.ufop.br/orientacoes-gerais-alunos-e-professores} - Informação Oficial com indicação das resoluções que tratam sobre cada assunto.
    \item Colegiado do curso - CC: \url{cocic@ufop.edu.br} e BIA: \url{cobia@ufop.edu.br}   
\end{enumerate}


\vspace{0.5cm}
\hrule
\vspace{0.3cm}

\noindent\footnotesize Fontes: Estatuto da UFOP (Resolução CUNI nº 1.868, de 17/02/2017, com alterações posteriores); Regimento Geral da UFOP (Resolução CUNI nº 1.959, de 28/11/2017, com alterações posteriores); página da PROGRAD --- Orientações Gerais (\url{https://www.prograd.ufop.br/orientacoes-gerais-alunos-e-professores}).

\end{document}